\chapter{Introduction}
\label{chap:introduction}

\section{Background Information}
\label{sec:background}

Multi-object tracking (MOT) in urban traffic scenes estimates the trajectories of vehicles across video frames while maintaining consistent object identities (IDs). In tracking-by-detection pipelines such as SORT \cite{bewley2016sort}, DeepSORT \cite{wojke2017deepsort}, ByteTrack \cite{zhang2022bytetrack}, and OC-SORT \cite{cao2023ocsort}, a detector provides per-frame bounding boxes and an association module links detections to active tracks. When a vehicle is occluded by another vehicle or roadside structure for $g \in [3, 90]$ consecutive frames, frame-by-frame IoU matching fails and the track enters a lost pool. Upon reappearance at confirmation frame $t_{\text{conf}}$, the tracker must solve a many-to-many re-association problem between $N$ lost tracks and $M$ unmatched candidate detections while simultaneously rejecting tracks that permanently exited the camera view and candidates corresponding to newly entered vehicles.

Linear Kalman filters with Mahalanobis gating struggle during medium-to-long occlusion gaps ($g \ge 16$ frames) for three reasons: (i) constant-velocity extrapolation accumulates nonlinear drift during braking, acceleration, or lane changes; (ii) uncalibrated process noise matrices mis-scale uncertainty ellipses across gap lengths; and (iii) unnormalized pairwise costs lack a differentiable mechanism to weigh competitive multi-track assignments against track/candidate dustbins (abstention).

\section{Motivation and Objective}
\label{sec:motivation_objective}

\textbf{Overall objective.} This capstone project designs, implements, and rigorously evaluates \textbf{GapLink-Lite}, a compact, self-developed neural network ($327{,}979$ trainable parameters, $< 1\text{M}$ budget) for vehicle identity re-association across occlusion gaps on the real-world \textbf{UA-DETRAC} benchmark \cite{wen2020uadetrac} ($100$ traffic sequences, $138{,}045$ annotated frames, $8{,}251$ vehicle trajectories, $1{,}201{,}341$ vehicle bounding boxes, and $0$ pedestrian annotations).

\textbf{Two-phase evaluation.} The project is evaluated in two phases. \emph{Phase~1} trains and tests the model on ground-truth vehicle boxes with controlled occlusions, under a pre-registered protocol, and freezes the result. \emph{Phase~2} deploys the frozen model behind a live YOLOv8n detector in a security-camera console (GapLink-VMS) and measures, on natural occlusions, how it compares with classical re-association approaches.

\textbf{Design constraints:}
\begin{enumerate}
    \item \textbf{Non-human tracking domain:} The entire pipeline operates strictly on road vehicles in UA-DETRAC ($0$ human/pedestrian boxes across all $100$ sequences, verified in \texttt{data/stats.json}).
    \item \textbf{Self-developed core model from random initialization:} Every module of GapLink-Lite---a 2-layer GRU motion encoder with learned gap embedding, a permutation-invariant DeepSets neighbor context encoder \cite{zaheer2017deepsets}, a constant-velocity residual Gaussian trajectory predictor, a pairwise + track/candidate dustbin scorer inspired by SuperGlue \cite{sarlin2020superglue}, and a 20-iteration log-domain Sinkhorn optimal transport layer \cite{sinkhorn1967,cuturi2013sinkhorn}---is trained from random initialization (\texttt{seeds 20260417--20260421}) without any pretrained weights or third-party tracker codebases.
\end{enumerate}

\textbf{Pre-registered hypotheses} (locked in \texttt{docs/PREREG.md} on 15/11/2026 before any test evaluation; verdicts follow the pre-registered wording):
\begin{itemize}
    \item \textbf{H1 (Event-level Top-1 on the hard subset):} Top-1 accuracy of GapLink-Lite on hard events (\texttt{is\_hard}) exceeds $\max(\text{C-2},\text{C-3},\text{C-4})$ by at least $3.0$ percentage points. \textbf{Outcome: Supported.} Hard-subset Top-1 $93.18\%$ vs.\ $85.85\%$ for the strongest baseline (C-4): $\Delta=+\mathbf{7.33}$~pp, 95\% paired bootstrap CI over the $15$ test camera locations $[6.13, 8.77]$~pp, whose lower bound also clears the $3.0$~pp margin (seed~1 = frozen \texttt{release/final.pt}).
    \item \textbf{H2 (Calibrated uncertainty coverage):} After gap-dependent calibration, empirical 90\% ellipse coverage lies in $[85\%, 95\%]$ in all three gap bins ($3$--$15$, $16$--$45$, $46$--$90$), measured on held-out validation half~B. \textbf{Outcome: Not supported (marginal miss in the long-gap bin).} Validation: $87.71\%$ / $87.64\%$ / $\mathbf{95.14\%}$; test: $87.87\%$ / $90.26\%$ / $\mathbf{95.86\%}$. The uncalibrated head lies inside the band on test ($86.42\%$ / $85.46\%$ / $88.02\%$), but this is a post-hoc observation and does not change the verdict.
    \item \textbf{H3 (Tracker-level IRR and WLR):} Identity Recovery Rate (IRR) of GapLink-Lite exceeds both C-2 and C-3 in at least $2/3$ gap bins, with $\Delta\text{WLR} \le +0.5$~pp. \textbf{Outcome: Supported (3/3 bins, full and partial occlusion).} Full occlusion: IRR $\mathbf{96.20\%}$ vs.\ $85.17\%$ (C-2) and $84.96\%$ (C-3), WLR $0.21\%$ vs.\ $4.32\%$ / $10.25\%$; location-bootstrap CI of $\Delta$IRR vs.\ C-3 $[9.10, 14.35]$~pp. Partial occlusion: IRR $\mathbf{95.99\%}$, WLR $0.36\%$, IDF1 $99.02\%$.
    \item \textbf{P1 (Protocol check):} With re-association disabled (C-1), IRR must be $<5\%$. \textbf{Outcome: Passed} (IRR $0.00\%$ on all $1{,}922$ injected intervals).
\end{itemize}

\section{Members and Role Assignments}
\label{sec:members_roles}

This is an individual capstone project. Table~\ref{tab:team_roles} summarizes the student's responsibilities, and Table~\ref{tab:contributions} explicitly separates self-developed components from external data and libraries.

\begin{table}[htbp]
    \centering
    \caption{Project responsibilities and deliverables}
    \label{tab:team_roles}
    \begin{tabularx}{\textwidth}{l p{3.5cm} X}
        \toprule
        \textbf{No.} & \textbf{Student} & \textbf{Responsibilities and Deliverables} \\
        \midrule
        1 & \memberOne & Formulate the GapLink-Lite architecture and pre-registered hypotheses, oversee data cleaning and camera-cluster splitting on UA-DETRAC, execute 5-seed training and ablation experiments on GPU, audit calibration/tracker metrics, build the GapLink-VMS application with live YOLOv8n, and finalize the report and video demo. \\
        \bottomrule
    \end{tabularx}
\end{table}

\begin{table}[htbp]
    \centering
    \caption{Explicit separation between self-developed components and external resources}
    \label{tab:contributions}
    \begin{tabularx}{\textwidth}{p{4.0cm} p{2.8cm} X}
        \toprule
        \textbf{Component} & \textbf{Origin} & \textbf{Implementation Details} \\
        \midrule
        UA-DETRAC Dataset & External Benchmark & $100$ real traffic video sequences ($960\times 540$, 25 fps) \cite{wen2020uadetrac}; cleaned to separate $528{,}798$ static ignore-zone rows from $1{,}201{,}341$ vehicle boxes. \\
        Camera-Clustered Split & Self-Developed & 64$\times$36 median background extraction + agglomerative clustering ($t=0.65$) into $25$ viewpoint clusters ($47$ train / $13$ val / $40$ test). \\
        GapLink-Lite Network & Self-Developed & 2-layer GRU + DeepSets Context + Residual Gaussian Predictor + SuperGlue-style Pairwise/Dustbin Scorer + Log-Sinkhorn ($327{,}979$ random-init params). \\
        5-View Partial Histograms & Self-Developed & Precomputed 128-bin HSV histograms for full box and 4 directional partial crops so occluded pixels never leak at inference. \\
        2-Stage Online Tracker & Self-Developed & Stage-1 IoU Hungarian + Stage-2 GapLink-Lite / C-2 / C-3 / C-4 re-association ($a_{\max}=95$) with IRR, WLR, IDF1, MOTA, and IDSW evaluation. \\
        \bottomrule
    \end{tabularx}
\end{table}

\section{Schedule and Milestones}
\label{sec:schedule_milestones}

The project was executed between 15 and 21 November 2026 as a sequence of verification gates (Table~\ref{tab:milestones}); each gate had to pass before the next started.

\begin{table}[htbp]
    \centering
    \caption{Schedule and milestones (verification gates)}
    \label{tab:milestones}
    \begin{tabularx}{\textwidth}{p{1.6cm} p{2.6cm} X}
        \toprule
        \textbf{Date} & \textbf{Gate} & \textbf{Verified Deliverables} \\
        \midrule
        15/11 & G1 Data & $100$ UA-DETRAC sequences converted and cleaned ($1{,}201{,}341$ vehicle boxes), $25$ camera clusters, location-disjoint split, 5-view HSV histograms. \\
        15/11 & G2 Events & Shared event sampler; validation ($6{,}000$) and test ($12{,}000$) event sets fixed and hashed; $120{,}038$ training events. \\
        15/11 & G3 Model & GapLink-Lite implemented ($327{,}979$ parameters); unit tests and 64-event overfit check ($100\%$ Top-1). \\
        15--16/11 & G4 Training & 5 seeds of the full model and 3 ablations (20 GPU runs); calibration fitted on validation. \\
        16/11 & G5 Freeze & Thresholds selected on validation, checkpoint frozen (\texttt{844fc391\ldots}). \\
        16/11 & G6 Phase-1 test & Single locked test run: event level, online tracker, ablations, robustness. \\
        18--20/11 & G7 Application & GapLink-VMS console with live YOLOv8n; forbidden-route, dwell and target-lock modules. \\
        21/11 & G8 Phase-2 test & Detector-aware checkpoint and live operating points selected on validation; single live test run; results frozen. \\
        \bottomrule
    \end{tabularx}
\end{table}
